\documentclass[a4paper,12pt]{article}
\usepackage[T2A]{fontenc}
\usepackage[utf8]{inputenc}
\usepackage[russian]{babel}
\usepackage{amsmath,amssymb}
\usepackage[a4paper,left=20mm,right=20mm,top=20mm,bottom=22mm,headheight=25pt,headsep=9mm]{geometry}
\usepackage{microtype,booktabs,tabularx,enumitem,fancyhdr,setspace,xcolor,tikz}
\usetikzlibrary{arrows.meta,positioning,shapes.geometric}
\definecolor{accent}{HTML}{008080}
\definecolor{darkaccent}{HTML}{004D4D}
\definecolor{textgray}{HTML}{333333}
\usepackage[hidelinks]{hyperref}
\hypersetup{pdftitle={Арифметика остатков и сравнения по модулю},pdfauthor={Лёвин Артём Александрович}}
\setlength{\parindent}{0pt}
\setlength{\parskip}{5pt}
\onehalfspacing
\setlist[itemize]{leftmargin=*,itemsep=1pt,topsep=3pt}
\setlist[enumerate]{leftmargin=*,itemsep=5pt,topsep=3pt}
\pagestyle{fancy}
\fancyhf{}
\fancyhead[L]{\small Арифметика остатков}
\fancyhead[R]{\small Ксения Клыкова}
\fancyfoot[C]{\small\thepage}
\renewcommand{\headrulewidth}{0.4pt}
\renewcommand{\headrule}{\hbox to\headwidth{\color{accent}\leaders\hrule height \headrulewidth\hfill}}
\newcommand{\topic}[1]{\par\medskip\noindent{\Large\color{darkaccent}#1}\par\smallskip{\color{accent}\hrule height 0.45pt}\medskip}
\newcommand{\example}[1]{\par\medskip\noindent\textcolor{darkaccent}{\textbf{#1}}\par\smallskip}
\newcommand{\reminder}[1]{\par\smallskip\textcolor{darkaccent}{\textbf{#1}}\par}
\newcommand{\modu}[1]{\pmod{#1}}
\begin{document}
\color{textgray}
\begin{titlepage}
\centering
\vspace*{25mm}
{\Large\color{darkaccent} МАТЕМАТИКА\par}
\vspace{15mm}
{\huge Чек-лист\par}
\vspace{10mm}
{\LARGE\color{accent} Арифметика остатков\\ и сравнения по модулю\par}
\vspace{15mm}
{\large Олимпиадные методы решения задач на делимость\par}
\vfill
{\large Ксения Клыкова\par}
\vspace{9mm}
{\large 25 ноября 2026 г.\par}
\vspace{18mm}
{\small Лёвин Артём Александрович\\ эксклюзивно для Ксении Клыковой\par}
\vspace{17mm}
\begin{tikzpicture}
\draw[accent,line width=0.9pt] (0,0) .. controls (2,0.8) and (4,-0.8) .. (6,0) .. controls (8,0.8) and (10,-0.8) .. (12,0);
\end{tikzpicture}
\vspace{5mm}
\end{titlepage}
\setcounter{page}{1}
\topic{1. Остаток и сравнение по модулю}
Пусть $a\in\mathbb Z$, $m\in\mathbb N$, $m\ge 1$. При делении с остатком существует единственная пара целых чисел $q,r$, для которой
\[a=mq+r,\qquad 0\le r<m.\]
Здесь $q$ — неполное частное, $r$ — остаток. Запись $m\mid a$ означает, что $a$ делится на $m$ без остатка.

Два целых числа сравнимы по модулю $m$, если их остатки при делении на $m$ совпадают:
\[a\equiv b\modu{m}\quad\Longleftrightarrow\quad m\mid(a-b).\]
\example{Почему это верно?}
Если $a=mx+r$, $b=my+r$, то $a-b=m(x-y)$. Обратно: если $a-b=mk$ и $b=my+r$, то $a=m(y+k)+r$. Следовательно, остаток у чисел один и тот же.
\example{Удобный представитель остатка}
\[98=11\cdot9-1,\qquad 98\equiv-1\equiv10\modu{11}.\]
Число $-1$ удобно для вычислений. Стандартный остаток равен $10$ и принадлежит множеству $\{0,1,\ldots,m-1\}$.
\topic{2. Основные свойства сравнений}
Если $a\equiv b\modu{m}$ и $c\equiv d\modu{m}$, то
\begin{align*}
 a+c&\equiv b+d\modu{m},&a-c&\equiv b-d\modu{m},\\
 ac&\equiv bd\modu{m},&a^k&\equiv b^k\modu{m}\quad(k\in\mathbb N).
\end{align*}
Для суммы и разности достаточно сложить или вычесть кратные $m$ разности. Например, для произведения $ac-bd=c(a-b)+b(c-d)$; оба слагаемых делятся на $m$. Свойство степеней следует из многократного умножения.
\clearpage
\topic{3. Вычисление остатков без больших чисел}
Главный приём: заменить каждое число удобным сравнимым с ним числом по одному и тому же модулю. Затем выполнить действия и привести результат к стандартному остатку.
\example{Пример 1. Сумма и произведение}
Найдём остаток от деления $98\cdot57+125$ на $11$.
\[98\equiv-1,\qquad57\equiv2,\qquad125\equiv4\modu{11}.\]
\[98\cdot57+125\equiv(-1)\cdot2+4=2\modu{11}.\]
Остаток равен $2$. Если вычесть $2$ из исходного числа или прибавить $9$, получится число, кратное $11$.
\example{Пример 2. Повторное приведение остатка}
\[137\cdot206+75\equiv2\cdot(-1)+3=1\modu{9}.\]
\[37^3+7\equiv2^3+2=10\equiv0\modu{5}.\]
Таким образом, первое выражение даёт остаток $1$, второе делится на $5$.
\topic{4. Подстановка в многочлены}
Если $n\equiv r\modu{m}$, то для любого многочлена $P$ с целыми коэффициентами $P(n)\equiv P(r)\modu{m}$.
\example{Пример 3. Отрицательный представитель}
Пусть $n\equiv-2\modu{7}$. Тогда
\[n^2+5n+1\equiv(-2)^2+5(-2)+1=-5\equiv2\modu{7}.\]
Поскольку $-5=7\cdot(-1)+2$, остаток равен $2$.
\example{Пример 4. Многочлен третьей степени}
При $n\equiv-1\modu{7}$ имеем
\[n^3+2n^2-5n+6\equiv-1+2+5+6=12\equiv5\modu{7}.\]
\clearpage
\topic{5. Когда можно сократить сравнение?}
Из $ca\equiv cb\modu{m}$ следует
\[a\equiv b\modu{m/d},\qquad d=\gcd(c,m)=\text{НОД}(c,m).\]
Если $d=1$, модуль остаётся прежним. Если $d>1$, модуль уменьшается; при $m/d=1$ условие выполняется для любых целых $a,b$.
\example{Почему требуется менять модуль?}
Из $m\mid c(a-b)$, после записи $m=dm_1$, $c=dc_1$, получаем $m_1\mid c_1(a-b)$. Числа $m_1,c_1$ взаимно просты, поэтому $m_1\mid(a-b)$.
\example{Пример 5. Опасность сокращения}
\[2\cdot1\equiv2\cdot4\modu{6},\qquad 1\not\equiv4\modu{6}.\]
Но $\gcd(2,6)=2$, поэтому верно $1\equiv4\modu{3}$.
\topic{6. Доказательство делимости}
Для доказательства $m\mid N$ достаточно показать $N\equiv0\modu{m}$.

Если $u,v$ взаимно просты, то одновременная делимость на $u$ и $v$ равносильна делимости на $uv$. В частности,
\[15\mid N\quad\Longleftrightarrow\quad
\begin{cases}N\equiv0\modu{3},\\N\equiv0\modu{5}.\end{cases}\]
Для невзаимно простых делителей гарантируется делимость на их НОК. Например, $12$ делится на $2$ и $4$, но $8\nmid12$.
\example{Пример 6. Квадрат нечётного числа}
Пусть $n=2k+1$, $k\in\mathbb Z$. Тогда
\[n^2=4k^2+4k+1=4k(k+1)+1.\]
Среди соседних целых чисел $k$ и $k+1$ одно чётное. Значит, $8\mid4k(k+1)$ и $n^2\equiv1\modu{8}$.
\clearpage
\topic{7. Десятичная запись и последняя цифра}
Для неотрицательного целого числа $N$ последняя цифра совпадает с остатком от деления на $10$:
\[N=10q+r,\qquad r\in\{0,1,\ldots,9\}.\]
Например, $125=10\cdot12+5$, $12346=10\cdot1234+6$.
\example{Степени десяти}
\[10\equiv1\modu{9},\qquad10\equiv-1\modu{11}.\]
Для целого $k\ge0$:
\[10^k\equiv1\modu{9},\qquad10^k\equiv(-1)^k\modu{11}.\]
По модулю $11$ при чётном $k$ получается $1$, при нечётном — $-1$ (стандартный остаток $10$).
\topic{8. Какие цифры возможны у квадрата?}
Может ли квадрат целого числа оканчиваться на $7$? Пусть $n=10q+r$, где $0\le r\le9$. Тогда
\[n^2=100q^2+20qr+r^2\equiv r^2\modu{10}.\]
Достаточно проверить десять вариантов $r$. Это полный перебор всех возможных остатков.
\begin{center}
\renewcommand{\arraystretch}{1.35}
\begin{tabular}{l*{10}{r}}
\toprule
$r$ &0&1&2&3&4&5&6&7&8&9\\
\midrule
$r^2$ &0&1&4&9&16&25&36&49&64&81\\
Остаток &0&1&4&9&6&5&6&9&4&1\\
\bottomrule
\end{tabular}
\end{center}
Возможные последние цифры квадрата: $0,1,4,5,6,9$. Цифры $2,3,7,8$ невозможны. В частности, квадрат целого числа не может оканчиваться на $7$.
\reminder{Приём полного перебора}
Для доказательства невозможности остатка переберите все классы от $0$ до $m-1$ и вычислите остаток выражения в каждом случае. Отдельных числовых примеров для общего доказательства недостаточно.
\clearpage
\topic{9. Делимость выражений с большими степенями}
Особенно удобно использовать сравнения с противоположными числами. Если $a\equiv-b\modu{m}$, то
\[a^n\equiv(-b)^n\modu{m}.\]
\[(-b)^n=\begin{cases}b^n,& n\text{ чётное},\\-b^n,& n\text{ нечётное}.\end{cases}\]
\example{Пример 7. Сумма степеней}
Докажем, что $2^n+3^n$ делится на $5$ при каждом положительном нечётном целом $n$.

Поскольку $2-(-3)=5$, имеем $2\equiv-3\modu{5}$. Возведём обе части сравнения в степень $n$:
\[2^n\equiv(-3)^n=-3^n\modu{5}.\]
Знак минус сохраняется благодаря нечётности показателя. Поэтому
\[2^n+3^n\equiv-3^n+3^n=0\modu{5}.\]
Требуемая делимость доказана для всех положительных нечётных $n$. В частности, $2^{10^3+1}+3^{10^3+1}$ делится на $5$, поскольку $10^3+1=1001$.
\topic{10. Типичные ошибки}
\begin{itemize}
\item Использовать отрицательный или слишком большой представитель как окончательный остаток: $-5\equiv2\modu{7}$, остаток $2$.
\item Путать равенство и сравнение: $98\equiv-1\modu{11}$, хотя $98\ne-1$.
\item Сокращать общий множитель без проверки НОД с модулем.
\item Объединять делимость на невзаимно простые числа через их произведение.
\item Терять скобки и знак при возведении отрицательных чисел в степень; проверяйте чётность показателя.
\item Доказывать общий факт по нескольким частным примерам; для перебора необходимо охватить все возможные остатки.
\end{itemize}
\clearpage
\topic{Алгоритм: вычисление остатка или доказательство делимости}
\vfill
\begin{center}
\begin{tikzpicture}[node distance=9mm, font=\small,
term/.style={ellipse,draw=accent,minimum width=43mm,minimum height=10mm,align=center},
actionnode/.style={draw=accent,rounded corners=2pt,text width=83mm,minimum height=13mm,align=center,inner sep=5pt},
dec/.style={diamond,draw=accent,aspect=2.9,text width=43mm,align=center,inner sep=1pt},
branch/.style={draw=accent,rounded corners=2pt,text width=61mm,minimum height=18mm,align=center,inner sep=5pt},
arr/.style={-{Stealth[length=2.5mm]},draw=darkaccent,line width=.65pt}]
\node[term] (start) {Начало};
\node[actionnode,below=of start] (m) {Определите модуль $m$ по условию};
\node[actionnode,below=of m] (rep) {Замените числа и переменные удобными сравнимыми представителями};
\node[actionnode,below=of rep] (calc) {Выполните сложение, вычитание, умножение и возведение в степень};
\node[actionnode,below=of calc] (norm) {Приведите результат к остатку $r$, где $0\le r<m$};
\node[dec,below=of norm] (goal) {Требуется доказать делимость?};
\node[branch,below=13mm of goal,xshift=-40mm] (rem) {Запишите остаток $r$};
\node[branch,below=13mm of goal,xshift=40mm] (div) {Проверьте: $r=0$ — делится; $r\ne0$ — не делится};
\node[term,below=44mm of goal] (end) {Ответ};
\draw[arr] (start)--(m); \draw[arr] (m)--(rep); \draw[arr] (rep)--(calc); \draw[arr] (calc)--(norm); \draw[arr] (norm)--(goal);
\draw[arr] (goal.west) -| node[pos=.22,above]{нет} (rem.north);
\draw[arr] (goal.east) -| node[pos=.22,above]{да} (div.north);
\draw[arr] (rem.south) |- (end.west); \draw[arr] (div.south) |- (end.east);
\end{tikzpicture}
\end{center}
\vfill
\clearpage
\topic{Тренировочный блок}
Решайте с помощью сравнений по модулю. В задачах на делимость и невозможность остатка приведите доказательство. Все переменные в заданиях — целые числа.
\example{Базовый уровень}
\begin{enumerate}
\item Найдите стандартный остаток от деления $-38$ на $7$.
\item Найдите остаток от деления $109\cdot46+124$ на $11$.
\item Найдите остаток от деления $43^3+12$ на $5$.
\end{enumerate}
\example{Средний уровень}
\begin{enumerate}[start=4]
\item Известно, что $n\equiv-3\modu{7}$. Найдите остаток от деления $n^2+4n+2$ на $7$.
\item Известно, что $n\equiv-2\modu{9}$. Найдите остаток от деления $n^3-2n^2+5n+4$ на $9$.
\item Известно, что $6x\equiv18\modu{24}$. Запишите все возможные стандартные остатки $x$ по модулю $24$.
\item Найдите остатки числа $10^{2026}$ при делении на $9$ и на $11$.
\end{enumerate}
\example{Повышенный уровень}
\begin{enumerate}[start=8]
\item Может ли квадрат целого числа оканчиваться на $8$? Обоснуйте ответ перебором возможных остатков.
\item Докажите, что $15\mid(10^{2k+1}+5)$ для любого целого $k\ge0$, отдельно проверив делимость на $3$ и на $5$.
\item Докажите, что $4^n+7^n$ делится на $11$ при каждом положительном нечётном целом $n$. Примените результат к $n=10^4+1$.
\end{enumerate}
\clearpage
\topic{Ответы}
\begin{center}
\renewcommand{\arraystretch}{1.28}
\begin{tabularx}{\linewidth}{@{}rX@{}}
\toprule
№ & Ответ \\
\midrule
1 & $4$\\
2 & $1$\\
3 & $4$\\
4 & $6$\\
5 & $5$\\
6 & $3,7,11,15,19,23$\\
7 & $1$ и $1$ соответственно\\
8 & Нет; возможны только цифры $0,1,4,5,6,9$.\\
9 & Делится на $15$ при всех указанных $k$.\\
10 & Делится на $11$ при всех указанных $n$, в том числе при $n=10001$.\\
\bottomrule
\end{tabularx}
\end{center}
\vspace{6mm}
\topic{Итоговая памятка}
\begin{center}
\renewcommand{\arraystretch}{1.3}
\begin{tabularx}{\linewidth}{@{}p{.34\linewidth}X@{}}
\toprule
Сигнал в задаче & Полезное действие\\
\midrule
Большие числа & Замените их удобными представителями по модулю.\\
Отрицательный результат & Приведите к диапазону от $0$ до $m-1$.\\
Многочлен & Подставьте представитель остатка переменной.\\
Общий множитель & Проверьте НОД множителя и модуля.\\
Последняя цифра & Используйте модуль $10$.\\
Доказательство делимости & Получите нулевой остаток.\\
Сумма больших степеней & Ищите противоположные представители, учитывайте чётность.\\
\bottomrule
\end{tabularx}
\end{center}
\end{document}
